\subsection{Online selection within a finite message chain}
\label{sec:chain-prefix}

The same recovered direction supports selection after earlier messages
have been processed. We retain the joint root, confirmation and traffic
interfaces of Section~\ref{sec:joint-continuation} and give the selected
direction a finite chain prefix. Earlier message keys may be disclosed;
the selected message's successor seed may be disclosed after its
derivation. The comparison tracks the seeds connecting those two
observation boundaries.

\paragraph{Prefix interface and selection.}
At the first allocation of a fresh traffic information string, designate
its material identity and a direction as before. Fix a deterministic
lifetime bound $1\le L\le 2^{64}-1$. With $s_0$ the designated traffic
seed, the prefix has the native equations
\begin{equation}
\begin{split}
 I_i&=D_M\Vert\texttt{chain}\Vert\operatorname{BE64}(i),\\
 P_i&=F(0,s_i),\qquad
 (s_{i+1},K_i)=\operatorname{Expand}(P_i,I_i,64),
 \quad 0\le i<L.
\end{split}
\label{eq:prefix-chain}
\end{equation}
Complete material identity, direction and index identify one original
step. Matching endpoint copies, failed computations and resumptions
retrieve that step's original result. Distinct traffic materials keep
their own histories, including when their epoch and index agree.

Immediately before each index's first private step computation, the
controller either continues the prefix or marks that index. This choice
uses its available history, including earlier disclosed message keys,
messages, authentication outcomes and counted public-oracle computation.
The mark is permanent and occurs at most once. Every computation at an
alias participates in the same first-computation ledger; this includes
transaction-private receive-side advancement through skipped indices.
An index already computed at any such alias has passed its selection point.

Before a mark, the seeds $s_i$ and Extract values $P_i$ are private
handles with only the uses in~\eqref{eq:prefix-chain}. The message keys
of steps passed without a mark may be returned and used arbitrarily
within the finite counted computation. At a marked step $\ell$, $K_\ell$
has the single-original-encryption, classical-decryption and exact-retry
interface of Corollary~\ref{cor:joint-first-message}. Its paired successor
$s_{\ell+1}$ becomes an available value, from which subsequent real chain
computation and disclosure may proceed. Private persistence uses the
same justified opaque interface. If the execution ends without a mark,
the selected application event and score are zero.

This timing separates observing earlier keys from choosing an already
observed target. It also fixes when each protected seed can first be
installed in a reduction. The root's other traffic entries, the opposite
direction, real confirmation and ACK checks, and successor-root disclosure
remain available as in the preceding joint continuation.

Let $\Delta_{\rm traffic}$ denote the first three terms of
\eqref{eq:joint-schedule-error}, capped at one, with $q_R,q_0,q_1$
bounding the expanded controller's entire execution. They compare the joint
root uses and the at most two traffic/confirmation PRK rows. For each
$i<L$, let $q_{s,i}$ and $q_{M,i}$ be deterministic upper bounds on
all ordinary compression queries in the corresponding complete
column and row reductions. In particular, these budgets include the
other real chain steps, known earlier keys, verification, subsequent
chain computation, preparation and terminal testing. Define
\begin{equation}
\begin{split}
 \Delta_{\rm prefix}(L)=\min\biggl\{1,\;&\Delta_{\rm traffic}
  +\sum_{i=0}^{L-1}\Bigl(
       \varepsilon_{\rm Hcol}(q_{s,i})\\
 &\hspace{30mm}+\varepsilon_{\rm Hrow}(2,q_{M,i},89)
                         \Bigr)\biggr\}.
\end{split}
\label{eq:prefix-error}
\end{equation}
The unused second traffic-row term is omitted under the same
one-class condition as in Theorem~\ref{thm:joint-message-leaf}.

\begin{theorem}[Online selected chain leaf]
\label{thm:online-chain-leaf}
For the prefix interface above, compare the real uniform-root schedule
with the experiment that replaces the joint root and traffic rows by
their preceding independent tables and gives each prefix step through
the mark two independent cached blocks $(s_{i+1},K_i)$. Computation
strictly after the marked step is real and starts from its available
successor seed. Every available terminal bit has probability gap at
most $\Delta_{\rm prefix}(L)$ between these experiments. Every common
available score in $[-1,1]$ has expectation gap at most
$2\Delta_{\rm prefix}(L)$.
\end{theorem}
\begin{proof}
The first three comparisons of
Theorem~\ref{thm:joint-message-leaf} retain the new controller as their
complete continuation. Their replies suffice to compute every chain
locally, including keys returned before the mark. They therefore cost
$\Delta_{\rm traffic}$. Apply
Lemma~\ref{lem:joint-expand-blocks} exactly to the traffic Expand table.
The designated $s_0$ is an independent private block; its parent row
has no remaining internal input containing that block.

Proceed in fixed index order. Before comparison $i$, every preceding
step reached before or at the mark has its independent block-pair
representation. If an earlier index was marked, step $i$ is already
in the real suffix and this comparison leaves that branch unchanged.
Otherwise the preceding block table, or the traffic table for $i=0$,
allocates a private uniform $s_i$. At that allocation the controller
has already decided whether the predecessor was marked. Hence the
column reduction can install its presampled hidden seed exactly on
the branch that continues the private prefix. On other branches it
uses locally sampled values. It never needs to predict the later mark.

Until step $i$ is first computed, its seed has only an opaque handle.
All available observations, including $K_{i-1}$ when $i>0$, can be
generated without its bytes. At the first computation, give $F(0,s_i)$
to the single-column challenger. Repeated computations use the same
answer. Its returned $P_i$ lets the reduction perform the real Expand,
all later steps, marking decisions and the complete public continuation.
If the execution never reaches this call, the challenge remains unused.
Thus this comparison costs $\varepsilon_{\rm Hcol}(q_{s,i})$ without
conditioning a primitive bound on either reachability or selection.

The fresh independent $P_i$ is then installed as the hidden key of
the row comparison. Its two Expand messages have lengths 57 and 89
bytes for every permitted index. The first reply supplies the input
prefix of the second query, so the reduction can generate both outputs
and the subsequent real continuation. This costs
$\varepsilon_{\rm Hrow}(2,q_{M,i},89)$. Exact block-table normalization
then removes the parent-row input containing $s_{i+1}$, giving an
independent pair. If step $i$ was marked, its successor is locally
available and all later steps remain real. Otherwise that successor
is the private handle for comparison $i+1$.

This construction includes skipped-key computation and failed private
attempts: their first step results are recorded before later publication
or commitment. All aliases reuse the same table entry. Independent
samples may coincide numerically; the comparisons preserve known-key
computations through the common public oracle and charge their queries.
No distinct-value conditioning is introduced.

The $2L$ comparisons use fixed worst-case budgets and recompute the
mark and terminal test within each game. Their triangle inequality,
together with the initial three comparisons, gives the bit bound.
The randomized test with success probability $(Z+1)/2$ gives the
stated score bound.
\end{proof}

\begin{corollary}[Selected message confidentiality after prefix disclosure]
\label{cor:online-chain-message}
Use the private-application and public-Resume boundary of
Corollary~\ref{cor:joint-first-message}, and let $A$ be an available
event independent of direct access to the hidden challenge bit $b$.
Include existence of a mark in $A$, and put
$Z=\mathbf1_A(2\mathbf1[g=b]-1)$. For the declared one-key,
one-encryption AEAD experiment,
\[
 |\mathbb E Z|\le
  2\Delta_{\rm prefix}(L)+2\varepsilon_{\rm AEAD}.
\]
For fresh authenticated installation under that marked message key,
\[
 \Pr[\mathsf{Bad}_{\rm selected\ msg}]
 \le\min\{1,\Delta_{\rm prefix}(L)+\varepsilon_{\rm INT}\}.
\]
Both primitive budgets include the entire adaptive selection and
subsequent continuation.
\end{corollary}
\begin{proof}
In the independent-pair experiment, a step is marked before its pair
is first generated. All earlier keys and marking decisions are
therefore available without $K_\ell$. Install the AEAD challenger's
key at that pair's message-key handle, sample its successor locally,
and compute the real suffix from that successor. Previous message
keys, confirmation keys and ACK keys are locally known. Exact retries
reuse the original encryption response recorded before its durable
outcome. The selected-key decryption interface supplies the same
authenticated marker or real error/plaintext handling as in
Corollary~\ref{cor:joint-first-message}.

For confidentiality, return $g$ on $A$ and a fair guess otherwise;
the centered success is the ideal experiment's $\mathbb E Z$.
For integrity, a new selected-key inbox installation absent from
the encryption ledger supplies the integrity experiment's winning
decryption pair. Apply Theorem~\ref{thm:online-chain-leaf} to the
respective score and Boolean event. This yields the two bounds.
\end{proof}

For $L=1$, the comparison has the same five primitive terms as
the selected first-message theorem. A fixed target index $\ell$
can use $L=\ell+1$; online selection uses a lifetime bound fixed
before execution. The native chain source uses the stated 56-byte
information strings, constant-width index encoding and zero AEAD
nonce, advances a send index once, and reuses exact outbox bytes on
Resume. Transaction-private receive advancement can compute skipped
keys before commitment; failed-transaction resumption retains their
first-computation identity in the proof ledger. The established native
next-rekey results retain their first-message controller; the
present theorem supplies the finite-chain continuation with its
explicit selection and observation interface.
